We introduce \emph{plane dihedrals}, a geometric visualization for two-dimensional spinors (qubit
states) embedded in $\mathbb{R}^3$.
Whereas the Bloch-sphere map identifies states that differ by a global $U(1)$ phase, plane dihedrals
make both the \emph{relative} phase between spinor components and a \emph{common} (fiber) phase
directly readable from geometry, and render the spinor’s characteristic $4\pi$ periodicity as a
visible double cover.
The representation consists of two offset circular sections on a sphere
(the \emph{plane-dihedral faces}): their radii encode the component moduli, while face markers encode
relative and common phase via rigid rotations.
We derive the construction from the $SU(2)$ formalism, relate it to the Bloch vector and the Hopf
fibration, and give an explicit forward and inverse mapping together with a complete rendering
pipeline.
We demonstrate static and animated views, including 2D projections that preserve the Born-rule
interpretation of the component weights.
An open-source Python reference implementation with interactive Jupyter notebooks is available at
\url{https://github.com/al-rigazzi/spinor_viz}.
